\section{Power as Causality}

In his important work \textit{The World as Power}, \textbf{John Woodroffe} lays out the philosophical principles of \textbf{Tantra}, while relating them to the Orthodox Hindu schools and even Western thought. Creation is the manifestation of Consciousness-Power and liberation is participation therein. In successive chapters, he describes Power as \textbf{Reality}, \textbf{Life}, \textbf{Mind}, \textbf{Matter}, \textbf{Causality}, \textbf{Continuity}, and \textbf{Consciousness}. Here we are interested in Causality , which is a special case of Continuity.

The first thing to take notice of is that, as \textbf{David Hume} so convincingly demonstrated, causality cannot be derived from purely empirical considerations. Thus the Enlightenment project of deriving all knowledge from observation of the physical world is impossible. Kant's odd solution was to situate causality in the mind, apart from the objects. Metaphysically, instead, we understand causality to be part of the cosmic order. Hence, we have in Causality an example of a real category, not derived by abstraction from empirical data.

The method of metaphysics is based on observation of states of consciousness. Unless certain states have been realized, its doctrines will appear opaque or speculative. Yet the refusal to make the efforts to achieve such states is as irrational as the bishop who refused to look through Galileo's telescope to observe the moons of Jupiter. Woodroffe's discussion of causality and continuity is too intricate for a single post, so we will have to be content to start with a few main points. Woodroffe summarizes the Vedantic position:

\begin{quotex}
Reality is Experience. Experience is a Universe. This Universe lives, moves and has its being in Consciousness or Cit. Cit therefore is Reality and the foundation of Reality. There is no inscrutable thing-in-itself beyond or behind Consciousness. Far from being unknowable, Reality is Cognition itself. 

\end{quotex}
Hence, the cosmic order as the objective arrangements among things is identical to the subjective representations of those arrangements. That is, Subject and Object are one; Knowing (cognition) and Being (reality) are identical. This is non-duality.

We begin by observing the causal relationship between things or events in consciousness, so we must accept the objective existence of causality. In our own volitional activity, we see directly the principle of causality; otherwise, we could not understand the motives for action. We observe the correspondence between the rational and practical activity of the self (knowledge and action).

Further, we can observe the sense of power within ourselves, that is, our power to effect change. So our understanding of causation is derived from our own conscious activity. This is still dualistic, since there is Mind (purusha) acting on Matter (prakriti). But Mind and Matter are not opposed, as they arise together in consciousness; rather, they form a scale. What is experienced, then, as matter is what is resistant to my Power, that is, a privation.

A deeper understanding is that the phenomenal universe is reducible to one Reality: the \textbf{Absolute} or \textbf{Brahman}. The world has no reality independent of Consciousness. This is not an abstraction of thought nor derived from experience. Woodroffe explains:

\begin{quotex}
We say that there is in fact only reality namely Consciousness, a position for which we may argue but which cannot be establilshed with certitude except by an actual or direct experience of unity. Those who seek to establish supersensible truths on any other ground must fail; just as those who argue against the validity of the individual's experience must fail.

\end{quotex}
We can summarize the understanding of causality from four perspectives:

\begin{enumerate}
\item In Material causation, the motion of one thing is communicated to another. 
\item In Efficient causation, Consciousness is effective and the Psycho-physical Principle operates on its own laws spontaneously. 
\item In the Vedanta, the efficient and material causes are understood as aspects of one Reality. 
\item From the absolute standpoint, there is no causation at all, since it transcends empirical reality. 
\end{enumerate}


\flrightit{Posted on 2011-04-20 by Cologero }

\begin{center}* * *\end{center}

\begin{footnotesize}\begin{sffamily}



\texttt{Izak on 2011-04-24 at 00:34 said: }

What do you recommend as being the best starting point with Woodroffe? Is The World as Power a good place to start? I am considering ordering The Serpent Power and Sakti \& Sakta alongside some sort of scholarly glossary on Hindu terminology.

By the way, I like this website — I got the recommendation to visit here from the comment board at Taki's Magazine.


\hfill

\texttt{Cologero on 2011-04-24 at 10:44 said: }

I would start with World as Power since it is more accessible to a Westerner without too much terminology; it exposes the metaphysical viewpoint of Tantra. In conjunction with Guenon's book on Hindu doctrine, it will provide a good foundation. Of course, Evola's Yoga of Power is a good synopsis; everything Evola knows about Tantra was taken from Woodroffe (at least the Hindu Tantras). Principles of Tantra is more systematic than the World as Power. Serpent Power assumes you are actively working with the Tantras with a qualified teacher.


\end{sffamily}\end{footnotesize}
